\documentclass[a4paper,11pt]{article}
\usepackage[utf8]{inputenc}
\usepackage[spanish]{babel}
\usepackage{geometry}
\usepackage{enumitem}
\usepackage{sectsty}
\usepackage{titlesec}
\usepackage{tikz}
\usepackage{array}

\geometry{
    a4paper,
    left=20mm,
    top=18mm,
    right=20mm,
    bottom=20mm
}

% Sin números de página
\pagestyle{empty}

% Estilo de secciones con números en círculos
\newcommand{\circled}[1]{%
    \tikz[baseline=(char.base)]{%
        \node[shape=circle,draw,inner sep=0.5pt,fill=black,text=white,font=\bfseries\footnotesize] (char) {#1};%
    }%
}

% Formato de secciones
\titleformat{\section}
{\normalfont\Large\bfseries}
{\circled{\thesection}}{0.5em}{}

\titlespacing*{\section}{0pt}{12pt}{8pt}

\titleformat{\subsection}
{\normalfont\normalsize\bfseries}
{}{0em}{}

\titlespacing*{\subsection}{8pt}{4pt}{4pt}

\begin{document}

% Encabezado del curso
\begin{center}
\small{Universidad de Granada}

\vspace{0.3cm}

\textbf{\large{Producción e Interacción Oral - Nivel 7}}

\vspace{0.4cm}

\small
\begin{tabular}{@{}ll@{}}
\textbf{Grupos:} & A (Aula 24) y B (Aula 08) \\
\textbf{Centro:} & Centro de Lenguas Modernas / Edificio A \\
\textbf{Profesor:} & Javier Benítez Lainez \\
\textbf{Tutorías:} & Martes 9:00-10:00 y 14:30-15:30 \\
& Sala de profesores / Email: benitezl@go.ugr.es
\end{tabular}
\end{center}

\vspace{0.5cm}
\section*{Ejercicios de Hipótesis en Pasado}

\vspace{0.4cm}

\section{Práctica de Pluscuamperfecto Subjuntivo y Condicional Compuesto C1}

\vspace{0.2cm}

---

\vspace{0.2cm}

\section{INTRODUCCIÓN}

\vspace{0.2cm}

Cuaderno de ejercicios para dominar la expresión de hipótesis en pasado usando pluscuamperfecto de subjuntivo y condicional compuesto.

\vspace{0.2cm}

---

\vspace{0.2cm}

\section{EJERCICIO 1: Transformación Básica}

\vspace{0.2cm}

Transforma situaciones reales en hipótesis pasadas:

\vspace{0.2cm}

\textbf{Modelo:} "No lo supe, no vine" → "Si \textbf{hubiera sabido}, \textbf{habría venido}."

\vspace{0.2cm}

1. "No estudié, no aprobé."

\textbf{Respuesta:} "Si \textbf{hubiera/hubiese estudiado}, \textbf{habría aprobado}."

\vspace{0.2cm}

2. "No teníamos dinero, no viajamos."

\textbf{Respuesta:} "Si \textbf{hubiéramos/hubiésemos tenido} dinero, \textbf{habríamos viajado}."

\vspace{0.2cm}

3. "No me avisaron, no fui."

\textbf{Respuesta:} "Si \textbf{me hubieran/hubiesen avisado}, \textbf{habría ido}."

\vspace{0.2cm}

4. "No llegó el tren, perdimos la conexión."

\textbf{Respuesta:} "Si \textbf{hubiera/hubiese llegado} el tren, \textbf{no habríamos perdido} la conexión."

\vspace{0.2cm}

---

\vspace{0.2cm}

\section{EJERCICIO 2: Selección de Tiempo Verbal}

\vspace{0.2cm}

Elige la forma correcta:

\vspace{0.2cm}

1. "Si ____ (había/hubiera) llegado antes, ____ (habría/había) visto el espectáculo."

\textbf{Respuesta:} hubiera, habría

\vspace{0.2cm}

2. "Ojalá ____ (había/hubiera) ____ (venido/venido) con nosotros."

\textbf{Respuesta:} hubiera/hubiese, venido

\vspace{0.2cm}

3. "Como si ____ (había/hubiera) ____ (pasado/pasado) nada."

\textbf{Respuesta:} hubiera/hubiese, pasado

\vspace{0.2cm}

4. "Si ____ (habríamos/hubiera) sabido, ____ (habríamos/habíamos) actuado diferente."

\textbf{Respuesta:} hubiéramos/hubiésemos, habríamos

\vspace{0.2cm}

---

\vspace{0.2cm}

\section{EJERCICIO 3: Historias de "Lo que Podría Haber Sido"}

\vspace{0.2cm}

Completa las historias con pluscuamperfecto de subjuntivo:

\vspace{0.2cm}

\textbf{Historia 1:}

"No ____ (obtener - yo) el trabajo. Si ____ (tener - yo) más experiencia, ____ (conseguir - yo) el puesto. De ____ (prepararme - yo) mejor para la entrevista, todo ____ (ser) diferente."

\vspace{0.2cm}

\textbf{Respuesta:} obtuve, hubiera/hubiese tenido, habría/hubiese conseguido; haberme preparado, habría/hubiese sido

\vspace{0.2cm}

\textbf{Historia 2:}

"Perdimos el vuelo. Si ____ (salir - nosotros) más temprano de casa, no ____ (perder - nosotros) el avión. De ____ (consultar - nosotros) el tráfico, ____ (llegar - nosotros) a tiempo."

\vspace{0.2cm}

\textbf{Respuesta:} hubiéramos/hubiésemos salido, habríamos/hubiésemos perdido; haber consultado, habríamos/hubiésemos llegado

\vspace{0.2cm}

---

\vspace{0.2cm}

\section{EJERCICIO 4: Arrepentimientos y Lamentos}

\vspace{0.2cm}

Usa "ojalá" + pluscuamperfecto de subjuntivo:

\vspace{0.2cm}

1. No viniste a la fiesta.

\textbf{Respuesta:} "¡Ojalá \textbf{hubieras/hubieses venido} a la fiesta!"

\vspace{0.2cm}

2. No me informé bien antes.

\textbf{Respuesta:} "¡Ojalá \textbf{me hubiera/hubiese informado} mejor!"

\vspace{0.2cm}

3. No tomamos esa decisión.

\textbf{Respuesta:} "¡Ojalá \textbf{hubiéramos/hubiésemos tomado} esa decisión!"

\vspace{0.2cm}

4. No te dije la verdad.

\textbf{Respuesta:} "¡Ojalá \textbf{te hubiera/hubiese dicho} la verdad!"

\vspace{0.2cm}

---

\vspace{0.2cm}

\section{EJERCICIO 5: Condicionales Mixtos}

\vspace{0.2cm}

Usa condicionales que combinan pasado y presente:

\vspace{0.2cm}

1. "Si ____ (estudiar - pasado) más cuando era joven, ____ (ser - presente) más exitoso ahora."

\textbf{Respuesta:} hubiera/hubiese estudiado, sería

\vspace{0.2cm}

2. "Si ____ (ser - presente) más responsable, ____ (terminar - pasado) el proyecto a tiempo."

\textbf{Respuesta:} fuera/fuese, habría/hubiese terminado

\vspace{0.2cm}

3. "Si no ____ (gastar - pasado) tanto dinero, ____ (tener - presente) más ahorros ahora."

\textbf{Respuesta:} hubiera/hubiese gastado, tendría

\vspace{0.2cm}

---

\vspace{0.2cm}

\section{EJERCICIO 6: Comparaciones con "Como si"}

\vspace{0.2cm}

Completa con pluscuamperfecto de subjuntivo:

\vspace{0.2cm}

1. "Actuó ____ ____ (ser) un experto."

\textbf{Respuesta:} como si \textbf{fuera/fuese}

\vspace{0.2cm}

2. "Se comportó ____ no ____ (pasar) nada."

\textbf{Respuesta:} como si \textbf{hubiera/hubiese pasado}

\vspace{0.2cm}

3. "Lo miró ____ no ____ (conocer) de nada."

\textbf{Respuesta:} como si \textbf{lo hubiera/hubiese conocido}

\vspace{0.2cm}

4. "Hablaba ____ ____ (tener) toda la razón."

\textbf{Respuesta:} como si \textbf{tuviera/tuviese}

\vspace{0.2cm}

---

\vspace{0.2cm}

\section{EJERCICIO 7: Narraciones de Alternativas}

\vspace{0.2cm}

Escribe una breve narración usando:

\begin{itemize}

  \item Mínimo 3 oraciones con pluscuamperfecto de subjuntivo

  \item Mínimo 1 "ojalá"

  \item Mínimo 1 inversión (de + haber + participio)

\end{itemize}

\vspace{0.2cm}

\textbf{Ejemplo:}

"Si \textbf{hubiera/hubiese tomado} ese camino, \textbf{habría llegado} más rápido. \textbf{Ojalá hubiera/hubiese sabido} esto antes. \textbf{De haberme informado}, \textbf{habría elegido} mejor. \textbf{Como si hubiera/hubiese sido} imposible anticiparlo."

\vspace{0.2cm}

---

\vspace{0.2cm}

\section{EJERCICIO 8: Role-Play de Arrepentimiento}

\vspace{0.2cm}

En parejas, uno expresa arrepentimiento y el otro muestra empatía:

\vspace{0.2cm}

\textbf{Persona A:} "Si ____ (hubiera/hubiese) ____ (hecho) X, ____ (habría/hubiese) ____ (pasado) Y."

\textbf{Persona B:} "Entiendo. ____ (Ojalá - tú) ____ (poder - perfecto) cambiar el pasado."

\vspace{0.2cm}

---

\vspace{0.2cm}

\section{EJERCICIO 9: Análisis de Errores}

\vspace{0.2cm}

Identifica y corrige los errores:

\vspace{0.2cm}

1. ❌ "Si \textbf{habría sabido}, habría venido."

✅ "Si \textbf{hubiera/hubiese sabido}, habría venido."

\vspace{0.2cm}

2. ❌ "Si \textbf{hubiera venido}, habría \textbf{venir} más temprano."

✅ "Si \textbf{hubiera/hubiese venido}, habría \textbf{venido} más temprano."

\vspace{0.2cm}

3. ❌ "Ojalá \textbf{había venido}."

✅ "Ojalá \textbf{hubiera/hubiese venido}."

\vspace{0.2cm}

---

\vspace{0.2cm}

\section{EJERCICIO 10: Conversaciones Espontáneas}

\vspace{0.2cm}

Responde estas preguntas con hipótesis pasadas:

\vspace{0.2cm}

1. ¿Qué habrías hecho de forma diferente en tu vida?

\textbf{Respuesta modelo:} "Si \textbf{hubiera/hubiese sabido} lo que sé ahora, \textbf{habría/hubiese tomado} diferentes decisiones."

\vspace{0.2cm}

2. ¿Qué pasaría si hubieras nacido en otro país?

\textbf{Respuesta modelo:} "Si \textbf{hubiera/hubiese nacido} en otro país, \textbf{habría/hubiese tenido} una vida completamente diferente."

\vspace{0.2cm}

3. ¿Qué habrías hecho si hubieras ganado la lotería el año pasado?

\textbf{Respuesta modelo:} "Si \textbf{hubiera/hubiese ganado} la lotería, \textbf{habría/hubiese viajado} por el mundo."

\vspace{0.2cm}

---

\vspace{0.2cm}

\textbf{Sesión 14: Ejercicios de Hipótesis en Pasado}

\textbf{Nivel C1 - Español Oral Avanzado}

\vspace{0.2cm}

\end{document}
